
\section{Feedback, Human Approval, and Runtime Governance}
Evaluation and approval are cross-cutting patterns. They can wrap chains, single agents, supervisors, or multi-agent systems.

\subsection{Evaluator--optimizer (reflection)}
A generator produces an artifact; tests and a rubric-based evaluator assess it; bounded revisions run when measurable improvement is feasible. A model judge is useful feedback but not an independent correctness or security oracle.

\begin{figure}[htbp]\centering
\resizebox{0.94\linewidth}{!}{%
\begin{tikzpicture}
\node[flowbox] (gen) at (0,0) {Generator\\Artifact vN};
\node[flowcontrol] (test) at (3.3,0) {Automated tests\\and evidence};
\node[flowbox] (judge) at (6.6,0) {Rubric evaluator};
\node[flowcontrol] (pass) at (6.6,-1.7) {Acceptance gate};
\node[flowbox] (rev) at (2.5,-1.7) {Bounded revision};
\draw[flowline] (gen)--(test);
\draw[flowline] (test)--(judge);
\draw[flowline] (judge)--node[flowlabel,right]{pass}(pass);
\draw[feedback] (judge)--node[flowlabel,above]{failed criterion}(rev);
\draw[flowline] (rev.west) .. controls (-0.8,-1.7) and (-0.8,0) .. (gen.west);
\end{tikzpicture}}
\caption{Evaluator--optimizer. A failed objective test cannot be overruled by evaluator prose.}
\label{fig:evaluator}
\end{figure}

Use objective tests whenever possible: schema conformance, compilability, unit tests, citations and evidence coverage. Terminate at a maximum revision count, non-improving score, expired deadline, or verified acceptance.

\subsection{Human-in-the-loop (HITL)}
HITL is a durable workflow state, not merely a prompt. Bind approval to the principal, exact operation, target resource, argument or patch digest, environment/commit version, and expiry. Any change invalidates the grant. Reauthorize immediately before execution to prevent time-of-check/time-of-use violations.

\begin{figure}[htbp]\centering
\resizebox{0.96\linewidth}{!}{%
\begin{tikzpicture}
\node[flowbox] (proposal) at (0,0) {Proposed sensitive action};
\node[flowcontrol] (pre) at (3.3,0) {Preflight policy and action digest};
\node[flowreview] (human) at (6.6,0) {Approve, deny, or expire};
\node[flowcontrol] (recheck) at (10,0) {Recheck ACL and version};
\node[flowbox] (audit) at (10,-1.7) {Execute and audit};
\draw[flowline] (proposal)--(pre);
\draw[flowline] (pre)--(human);
\draw[flowline] (human)--node[flowlabel,above]{approved}(recheck);
\draw[flowline] (recheck)--(audit);
\end{tikzpicture}}
\caption{Action-specific approval gate. Denial and expiry terminate or return to review.}
\label{fig:approval}
\end{figure}

\subsection{Reliability and failure semantics}
\begin{itemize}
\item Enforce wall time, model calls, token/cost budget, worker fan-out and iteration limits in software.
\item Retry only within a deadline; distinguish read-only, idempotent and non-idempotent tools. Writes require idempotency keys and reconciliation.
\item Represent malformed output, policy denial, tool failure, timeout, cancellation, stale approval and verification failure as distinct outcomes.
\item Persist checkpoints and make partial failures explicit. An incomplete join must not be treated as success.
\item Treat external content, tool messages, and agent memory as untrusted data.
\end{itemize}

\section{Interoperability: MCP, A2A, and the Harness}
MCP is a capability interface connecting application hosts and clients to tool/resource/prompt servers. The July 2026 specification documents a host--client--server architecture \cite{mcp}. A2A supports agent discovery and task collaboration across independent agent services \cite{a2a}. Neither chooses the orchestration graph or automatically grants permissions. Verify deployed protocol versions before implementing wire-level integrations.

\begin{figure}[htbp]\centering
\resizebox{0.94\linewidth}{!}{%
\begin{tikzpicture}
\node[flowcontrol,text width=33mm] (host) at (0,0) {Harness: orchestration and identity};
\node[flowbox] (client) at (3.8,1) {MCP client};
\node[flowbox] (server) at (7.1,1) {MCP tool/resource server};
\node[flowbox] (a2a) at (3.8,-1) {A2A client};
\node[flowbox] (peer) at (7.1,-1) {Remote agent service};
\draw[flowline] (host)--(client);
\draw[flowline] (client)--node[flowlabel,above]{MCP}(server);
\draw[flowline] (host)--(a2a);
\draw[flowline] (a2a)--node[flowlabel,above]{A2A}(peer);
\draw[feedback] (peer.south) |- (host.south);
\end{tikzpicture}}
\caption{Protocols expose capabilities and collaboration; the application owns policy.}
\label{fig:interoperability}
\end{figure}

\begin{table}[htbp]\small\centering
\begin{tabularx}{\linewidth}{@{}p{23mm}Y Y@{}}\toprule
\textbf{Component} & \textbf{Provided capability} & \textbf{What it does not imply}\\\midrule
Model API & Inference and structured tool proposals & Permission to execute tools\\
Agent harness & State machine, tool dispatch, approvals and tracing & Model-generated correctness\\
MCP & Tools/resources/prompts over a host--server boundary & Agent-to-agent delegation or blanket ACL\\
A2A & Agent discovery, messages, tasks and artifacts & Shared internals or inherited identity\\
AGENTS.md / skills & Repository instructions and reusable procedures & Trusted execution or elevated privileges\\\bottomrule
\end{tabularx}
\caption{Protocols, runtime, instructions and topology solve different problems.}
\end{table}

\section{Security Architecture and Threat Model}
OWASP describes prompt injection, tool abuse, memory poisoning, exfiltration, excessive agency, and cascading failures as practical agent risks \cite{owasp}. NIST AI RMF and the Generative AI Profile provide organizational risk-management context \cite{nist,nistgenai}.

\begin{table}[htbp]\small\centering
\begin{tabularx}{\linewidth}{@{}p{25mm}Y Y@{}}\toprule
\textbf{Threat} & \textbf{Typical exposure} & \textbf{Required defense}\\\midrule
Prompt injection & Web documents, tools, agent messages & Separate data/instructions, policy outside prompts\\
Excessive agency & ReAct, planners, supervisors & ACL, caps, timeout, sandbox and approval\\
Credential escalation & Workers and A2A peers & Per-worker identity, receiver-side authorization\\
Memory poisoning & Shared or persistent state & Provenance, isolation, write policy, expiry\\
Denial of wallet & Fan-out and reflection & Concurrency, turn, token and cost limits\\
Evidence laundering & Synthesis, voting and evaluation & Source-linked evidence and independent tests\\
Approval bypass & Remediation or deployment & Action-bound expiring approval and version recheck\\\bottomrule
\end{tabularx}
\caption{Threat-informed controls apply across orchestration choices.}
\end{table}

\subsection{Non-negotiable invariants}
\begin{enumerate}
\item A model proposes actions; a separate executor authorizes and executes them.
\item Authorization checks current identity, resource, action, policy and arguments.
\item Each worker gets isolated state, narrow capabilities and auditable identity.
\item Source content never overrides the instruction hierarchy or enterprise policy.
\item Every loop has a deadline, budget, cancellation and terminal error state.
\item High-impact changes require exact approval, revalidation and an audit trail.
\item Logs redact secrets and do not collect private model reasoning by default.
\end{enumerate}
